\documentclass{article}

% NeurIPS 2025 template (preprint mode for arXiv).
\usepackage[preprint]{neurips_2025}

% ---- Packages ----
\usepackage{amsmath,amssymb,amsfonts}
\usepackage{graphicx}
\usepackage{booktabs}
\usepackage{multirow}
\usepackage{xcolor}
\usepackage{hyperref}
\usepackage{url}
\usepackage{microtype}
\usepackage{caption}
\usepackage{subcaption}
\usepackage{enumitem}
\usepackage{algorithm}
\usepackage{algpseudocode}
\usepackage{tikz}

\hypersetup{colorlinks=true, linkcolor=blue, citecolor=blue, urlcolor=blue}

% ---- Convenience macros ----
\newcommand{\cart}{\textsc{CART}}
\newcommand{\tbd}[1]{\textcolor{gray}{\textit{TBD: #1}}}
\newcommand{\pplwiki}{PPL$_\text{wiki}$}
\newcommand{\ppltiny}{PPL$_\text{tiny}$}
\newcommand{\ppledu}{PPL$_\text{edu}$}

% ---- Title ----
\title{CART: Context-Anchored Recurrent Transformer\thanks{Code, training scripts, and the experiment database are available at \url{https://github.com/ccapps42/CART}.}\\
A Parameter-Efficient Architecture with Learned Stability}

\author{%
  Chad A. Capps\\
  Independent Research\\
  \texttt{Chad.Capps@me.com}
}

% ============================================================
\begin{document}
\maketitle

% ============================================================
\begin{abstract}
We present \cart{} (Context-Anchored Recurrent Transformer), a parameter-efficient language model architecture that reuses a single shared core block $R$ times across depth. Unlike prior looped transformers that recompute key-value representations at every loop iteration, \cart{} computes $K$ and $V$ once from a multi-layer prelude and reuses them throughout the recurrent core via multi-head latent attention (MLA) cross-attention. This design separates context encoding from iterative refinement, reduces per-loop FLOPs, and provides a stable attention anchor for all loop iterations. Stability of the recurrent computation is ensured by a learned Linear Time-Invariant (LTI) gate, which empirically settles on a narrow range of spectral radius values across all model scales \tbd{Stage 2 final values: $\varrho \in [0.79, 0.82]$ from completed configs; awaiting d=1024 to finalize. The Stage 1 ``universal $\varrho \approx 0.893$'' claim was a 3{,}000-step artifact---rho continues to decay through full training. Stage 2 also reveals weak R-dependence not visible at Stage 1.}

We conduct a two-stage sweep across $d_{\text{model}} \in \{256, 512, 768, 1024\}$, prelude depths $P \in \{6\}$, and loop counts $R \in \{6, 8, 10\}$, training all models on $\sim$1B tokens of mixed-domain text at sequence length 1024. Stage 1 (3{,}000 steps) validates \cart{}'s design axes and screens hyperparameters. Stage 2 (30{,}500 steps) provides the primary comparison against Dense transformer baselines.

Preliminary Stage 2 results show \cart{} outperforms density-matched Dense baselines at $d \geq 512$ on all three evaluation metrics, with the advantage most pronounced on natural-language text (\pplwiki{}: $-9\%$, \ppledu{}: $-16\%$ at $d=512$). Complete Stage 2 results across all scales are \tbd{pending training completion, est.\ May 2026}.

A key empirical finding is that $P$ (prelude depth) dominates $R$ (loop count) across all scales in Stage 1, with $R$ delivering increasing returns at larger $d$. The LTI gate discovers a memory timescale without supervision, and this timescale adapts to model complexity but not strongly to loop count, a novel emergent behavior we characterize across 64 configurations.
\end{abstract}

% ============================================================
\section{Introduction}
\label{sec:intro}

The dominant paradigm in language model architecture is the stacked transformer~\citep{vaswani2017attention}: $L$ unique layers, each with independent parameters, applied sequentially. This design scales predictably with $L$, but stored parameter count grows linearly with depth. At inference, every unique layer must be loaded from memory, making model size the primary bottleneck for deployment on memory-constrained hardware.

An alternative approach, depth-recurrent or looped transformers, reuses a single set of parameters across multiple passes through a shared block. The Universal Transformer~\citep{dehghani2018universal} pioneered this direction: a single block applied $R$ times to a token sequence. If the $R$-th application produces a representation of comparable quality to a $R$-layer unique model, then the recurrent model has achieved $R$-fold parameter compression at the cost of sequential compute.

The core challenge is making weight-sharing work. Naive looped transformers underperform their depth-matched counterparts~\citep{zeitoun2026hyperloop}, because a single block applied identically $R$ times cannot adapt its behavior across loop iterations. Recent solutions include hyper-connections~\citep{zhang2024hyperconnections,zeitoun2026hyperloop,schwethelm2026isodepth}, loop-index embeddings~\citep{su2024roformer,gomez2026openmythos}, and explicit stability constraints~\citep{prairie2026parcae,gomez2026openmythos}. The ``prelude / shared-weight core / coda'' topology used in this paper is also shared by several concurrent works~\citep{geiping2025latent,prairie2026parcae,gomez2026openmythos}; what distinguishes \cart{} is how the prelude representation enters the loop. All prior work in this family injects the prelude embedding $e$ via addition or concatenation into the running state, which then \emph{self-attends}. \cart{} instead computes $K, V$ once from $e$ and has the loop body \emph{cross-attend} to those frozen tensors.

\paragraph{Architectural contributions.}
\cart{} incorporates the proven mechanisms above (hyper-connections, loop-index embeddings, MLA compression) and contributes three architectural elements that, in combination, are not present in any prior or concurrent work we have surveyed:

\begin{enumerate}[leftmargin=*, topsep=2pt, itemsep=1pt]
  \item \textbf{Once-computed KV anchor with shared-weight cross-attending core.} The prelude computes $K, V$ once via MLA compression and holds them fixed across all $R$ core iterations. All surveyed concurrent looped architectures (Geiping~\citep{geiping2025latent}, Hyperloop~\citep{zeitoun2026hyperloop}, Parcae~\citep{prairie2026parcae}, OpenMythos~\citep{gomez2026openmythos}, and SpiralFormer~\citep{spiralformer2026}) recompute $K, V$ from the evolving hidden state at each loop step. The closest precedent for shared-weight iterated cross-attention is the Perceiver family~\citep{jaegle2021perceiver,jaegle2022perceiverio}, which re-projects $K, V$ from the input array at each iteration rather than caching the projected tensors and is non-autoregressive. \cart{}'s frozen anchor separates context encoding from iterative refinement, reduces per-loop FLOPs by eliminating redundant KV projections, and provides a stable attention target throughout recurrence.
  \item \textbf{Sigmoid-gated learned stability.} A sigmoid-gated LTI filter guarantees spectral radius $\varrho < 1$ without structural constraints: $h \leftarrow \sigma(a) \odot h_\text{in} + f_\text{transformer}(h)$. Parcae~\citep{prairie2026parcae} and OpenMythos~\citep{gomez2026openmythos} impose stability architecturally via zero-order-hold (ZOH) discretization of a state matrix; \cart{}'s gate instead \emph{learns} its convergence rate from data.
  \item \textbf{Causal cross-attention requirement.} In autoregressive recurrent models where prelude and core share the same token sequence, non-causal cross-attention leaks future tokens to the recurrent state. We document this constraint and its empirical consequences (Section~\ref{sec:arch_notes}).
\end{enumerate}

\paragraph{Empirical findings.}
This paper reports four findings about looped transformers that, to our knowledge, have not been demonstrated in prior work:

\begin{enumerate}[leftmargin=*, topsep=2pt, itemsep=1pt]
  \item \textbf{Universal learned memory timescale.} \tbd{Stage 2 results supersede the original Stage 1 claim: at full training budget, $\varrho$ settles in the narrower band $[0.79, 0.82]$ rather than at $\varrho \approx 0.893$, and shows weak dependence on $R$ that Stage 1's 3{,}000-step measurement missed. Awaiting d=1024 Stage 2 completion to finalize the corrected claim.} Across all 64 Stage 1 configurations (four model widths $d \in \{256, 512, 768, 1024\}$ with $R \in \{2, 4, 6, 8\}$ and $P \in \{2, 3, 4, 6\}$), the LTI gate's spectral radius converges to $\varrho \approx 0.893$, indistinguishable from the $R{=}6$ ``ideal'' value $0.5^{1/R} = 0.8909$. The convergence rate adapts gradually with model width ($+0.0042$ from $d{=}256$ to $d{=}1024$) but is essentially insensitive to loop count.
  \item \textbf{Prelude depth dominates loop count.} The hyperparameter ordering $P{=}6 > P{=}4 > P{=}3 > P{=}2$ holds without exception across all $R$ values and all four scales. The pattern, that representational depth in the encoder matters more than recurrence depth in the core, has not been quantified across scales in prior work.
  \item \textbf{The benefit of recurrence depth grows monotonically with model width.} \tbd{This claim is from Stage 1 (3{,}000-step, very sub-Chinchilla measurements). Stage 2 partial data (d $\leq$ 768 complete, 1B tokens, still sub-Chinchilla on effective params at d=768) shows R-benefit absent or slightly negative within $R \in \{6, 8, 10\}$, contradicting the Stage 1 prediction. Whether R-benefit reappears at Chinchilla-class training is the subject of a planned extended-training experiment; the Stage 1 claim and the Stage 2 reinterpretation will be reconciled once that experiment completes.} The improvement from $R{=}2$ to $R{=}8$ ranges from $-0.25\%$ at $d{=}256$ (slight regression) to $+5.24\%$ at $d{=}1024$. This predicts that the value of looping increases with scale, a direction that the iso-depth scaling literature~\citep{schwethelm2026isodepth} measures via $\phi$ but has not previously characterized as a function of model width.
  \item \textbf{Natural-language advantage at parameter parity.} At $d{=}512$ with all 6 complete Stage 2 seeds, \cart{} outperforms a parameter-matched Dense baseline by $\sim 9\%$ on Wikipedia perplexity and $\sim 16\%$ on FineWeb-Edu perplexity, with only $\sim 1.5\%$ improvement on the synthetic narrow-vocabulary TinyStories set. The natural-language-vs-synthetic gap (defined precisely in Section~\ref{sec:eval_terminology}) suggests fixed-point recurrence extracts representations that benefit harder language-modeling targets more than they benefit easier ones.
\end{enumerate}

This paper presents \cart{}'s design, a systematic Stage 1 sweep across 64 configurations at four scales, and preliminary Stage 2 results against matched Dense baselines. Complete Stage 2 results will update this preprint upon training completion.

% ============================================================
\section{Related Work}
\label{sec:related}

\paragraph{Universal and looped transformers.}
\citet{dehghani2018universal} introduced the Universal Transformer, applying a single block recurrently with adaptive computation time. Subsequent work explored stability, efficiency, and training dynamics. \citet{jeddi2026loopformer} introduced budget-conditioned looping with shortcut-consistency training. \citet{chen2026deeper} demonstrates that shared-weight blocks iterated in latent space exhibit sharp compositional generalization absent in standard transformers. \citet{han2026hierarchical} empirically validates that hierarchical recurrence (distinct boundary layers + looped core) outperforms flat uniform looping, which directly supports \cart{}'s prelude/coda design.

\paragraph{Geiping recurrent-depth latent reasoning.}
\citet{geiping2025latent} introduced the prelude/core/coda terminology adopted in this paper, scaling a recurrent-depth language model to 3.5B parameters with the layout $(\ell_p, \ell_r, \ell_c) = (2, 4, 2)$. Their core block is shared-weight and iterated, but the prelude embedding $e$ is reinjected at each iteration via concatenation through an adapter $A: \mathbb{R}^{2h} \to \mathbb{R}^h$ mapping $[s_{i-1}; e]$, after which standard causal self-attention recomputes $K, V$ from the resulting state. Stability is achieved via a sandwich-RMSNorm pattern rather than a learned gate. \cart{} adopts the same three-zone layout but replaces the concat-and-self-attend mechanism with a frozen-KV cross-attention anchor and a sigmoid LTI gate.

\paragraph{Perceiver-style architectures.}
\citet{jaegle2021perceiver,jaegle2022perceiverio} introduced shared-weight iterated cross-attention from a learned latent array to a fixed input array, the closest structural precedent for \cart{}'s shared-weight cross-attending core. Two differences are essential. First, Perceiver re-projects $K, V$ from the input at every cross-attention call (the projection matrices are weight-tied; the resulting tensors are not cached). \cart{} caches the projected $K, V$ tensors themselves. Second, Perceiver's ``input array'' is a separate modality (bytes, images) attended to by a learned latent set, whereas \cart{}'s prelude output is itself a transformer encoding of the same token sequence used by the recurrent core, and the architecture is fully autoregressive and causal.

\paragraph{Hyperloop Transformers.}
\citet{zeitoun2026hyperloop} propose a structurally similar three-zone architecture (begin/middle/end, 25/50/25\% parameter split) with hyper-connections at loop boundaries. The critical difference from \cart{}: Hyperloop's middle block uses standard self-attention, recomputing $K, V$ from the current hidden state at every loop iteration. \cart{} computes $K, V$ once from the prelude and reuses them, reducing per-loop FLOPs and separating context encoding from iterative refinement. Additionally, \cart{} provides an explicit spectral radius guarantee via the LTI gate, which Hyperloop does not.

\paragraph{Parcae.}
\citet{prairie2026parcae} independently arrive at the looped prelude/coda structure and address stability via negative-diagonal parameterization of a state transition matrix, discretized via zero-order hold. This architecturally constrains $\varrho < 1$. \cart{} learns stability via a sigmoid gate: $h \leftarrow \sigma(a) \odot h_\text{in} + f_\text{transformer}(h)$, where $\varrho$ emerges from training rather than being structurally imposed. Parcae uses variable recurrence depth (Poisson-sampled) and truncated BPTT; \cart{} uses fixed $R$ and full backpropagation, which \citet{schwethelm2026isodepth} show preserves higher effective loop value ($\phi = 0.65$ vs.\ $0.38$ for truncated).

\paragraph{OpenMythos.}
\citet{gomez2026openmythos} is an open-source reconstruction of a hypothesized recurrent-depth architecture combining MLA, hyper-connections, loop-index embeddings, and LTI-stable injection. It is the most architecturally similar prior work to \cart{}. Two differences are architecturally significant. First, OpenMythos recomputes $K, V$ from the evolving hidden state at each loop iteration; \cart{} computes $K, V$ once from the prelude and holds them fixed. Second, OpenMythos implements LTI stability via ZOH discretization of a learned state matrix (the same mechanism as Parcae), whereas \cart{}'s sigmoid gate learns its convergence rate without structural constraints. \cart{} was informed by OpenMythos and shares its component vocabulary, but diverges in both the KV treatment and the stability mechanism.

\paragraph{Cross-layer parameter sharing without recurrent semantics.}
ALBERT~\citep{lan2020albert} shares parameters across stacked transformer layers as a compression strategy. The architecture is otherwise a standard non-recurrent encoder: there is no fixed-point semantics, no LTI gate, no prelude/coda separation, and no notion of ``iterating'' the shared block. \cart{} uses cross-layer weight sharing in the same mechanical sense (a single block is reused $R$ times) but interprets the reuse as recurrence with an explicit context anchor and a stability guarantee.

\paragraph{State-space models and linear-attention recurrence.}
Mamba~\citep{gu2023mamba}, Mamba-2~\citep{dao2024mamba2}, and RWKV~\citep{peng2023rwkv} achieve infinite-context reasoning via recurrence over the \emph{sequence} dimension, with parallel-scan training algorithms that recover transformer-like throughput. \cart{}'s recurrence is over the \emph{depth} dimension: $R$ shared-weight passes over a fixed-length token sequence. The two are complementary rather than competing: SSMs trade attention's quadratic context cost for fixed-state memory, while \cart{} trades unique-layer parameter cost for sequential depth iteration. A natural future direction is combining the two (an SSM-based prelude feeding a \cart{}-style recurrent core), but this paper restricts attention to a transformer prelude with MLA self-attention.

\paragraph{Other concurrent variants.}
Three additional concurrent works share aspects of \cart{}'s design but diverge on the key mechanistic details. \citet{spiralformer2026} adopts the same three-zone topology with shared-weight loop and pre-/post-loop blocks, but uses multi-resolution recursion with self-attention on the evolving state. \citet{plt2025parallel} reuses $K, V$ across all loop iterations (the closest precedent for the once-computed-KV idea), but the $K, V$ are projected from the first iteration of the loop block itself, not from a unique prelude over the same tokens, and serve a sliding-window/global hybrid attention pattern rather than a recurrent state anchor. \citet{rdvit2026} applies LTI-stable state injection to vision (semantic segmentation) using a single shared block with optional MoE; the cross-modal nature and dense-prediction setting make a direct comparison difficult, but the LTI-stable framing parallels \cart{}'s. None of these works combine all of: a unique-layer prelude producing a frozen KV anchor, a shared-weight cross-attending core, and a learned sigmoid LTI gate.

\paragraph{Iso-depth scaling laws.}
\citet{schwethelm2026isodepth} measure the recurrence-equivalence exponent $\phi$, defined as how much one additional loop contributes relative to a unique layer. They find $\phi = 0.46$ for vanilla looped models, rising to $\phi = 0.65$ with hyper-connections. \cart{} incorporates hyper-connections, and we plan to compute $\phi$ for \cart{} using their framework once Stage 2 is complete (Section~\ref{sec:discussion}).

\paragraph{MLA compression.}
\citet{liu2024deepseek} introduced Multi-head Latent Attention for KV compression in mixture-of-experts models. \cart{} applies MLA to the prelude's KV projection ($d_{kv} = d/4$), making the once-computed context anchor cheaper to produce and store.

% ============================================================
\section{Architecture}
\label{sec:arch}

\begin{figure}[t]
  \centering
  \includegraphics[width=0.55\linewidth, height=0.55\textheight, keepaspectratio]{../figures/cart_architecture.png}
  \caption{\cart{} architecture. The prelude produces context anchor $e$ via $P$ unique layers. KV projections are computed once from $e$ and reused across all $R$ core iterations. The coda produces final token representations.}
  \label{fig:arch}
\end{figure}

\cart{} divides depth into three sequential zones (Figure~\ref{fig:arch}):

\[
\text{Embedding} \to \underbrace{P \text{ unique layers}}_{\text{Prelude}} \to \underbrace{\text{KV}(e)}_{\text{once}} \to \underbrace{R \times \text{CoreBlock}}_{\text{Recurrent Core}} \to \underbrace{1 \text{ layer}}_{\text{Coda}} \to \text{Logits}
\]

\subsection{Prelude}

The prelude consists of $P$ unique transformer layers, each with MLA self-attention and SwiGLU FFN~\citep{shazeer2020glu}. MLA compresses the key-value cache: instead of projecting $K, V \in \mathbb{R}^{T \times d}$ directly, it first projects $x \in \mathbb{R}^{T \times d}$ to a compressed latent $c \in \mathbb{R}^{T \times d_{kv}}$ (with $d_{kv} = d/4$), then expands to full-rank $K, V$ for attention. The query $Q$ is full-rank. This reduces the cost of the KV projection by $4\times$.

The final prelude output $e \in \mathbb{R}^{T \times d}$ serves as the context anchor for the recurrent core. Before entering the loop, $K$ and $V$ are computed from $e$ once via a single MLA KV Projection module and held fixed for all $R$ iterations.

\subsection{Recurrent Core}

A single \texttt{CoreBlock} is applied $R$ times with shared weights. At each loop step $\ell \in \{1, \ldots, R\}$:

\begin{enumerate}[leftmargin=*, topsep=2pt, itemsep=1pt]
  \item \textbf{HyperConnection blend.} A ring buffer maintains the last $n_h = 3$ hidden states $\{h_{\ell-1}, h_{\ell-2}, h_{\ell-3}\}$. Softmax-weighted residual combination produces a blended input $\tilde{h}$, initialized with residual weights $[1, 0, 0]$~\citep{zhang2024hyperconnections}.

  \item \textbf{Loop Index Embedding (LIE).} A sinusoidal embedding of the loop index $\ell$ is added to $\tilde{h}$, providing a loop-step signal analogous to RoPE applied over recurrence depth rather than sequence position~\citep{su2024roformer,gomez2026openmythos}.

  \item \textbf{MLA cross-attention.} $Q$ is projected from $\tilde{h}$; $K, V$ are the prelude-derived cache from step 1. Attention is computed with \texttt{is\_causal=True} (see Section~\ref{sec:arch_notes}).

  \item \textbf{SwiGLU FFN.} Hidden dimension $d_{ff} = \lceil (8/3) \cdot d / 256 \rceil \times 256$ following standard practice.

  \item \textbf{LTI gate.} The output is passed through:
  \[
    h_\ell = \sigma(a) \odot \tilde{h} + f_\text{attn+FFN}(\tilde{h})
  \]
  where $a \in \mathbb{R}^d$ is a learned gate. The sigmoid gating ensures $0 < \sigma(a) < 1$, which, combined with the transformer output scaling, guarantees the spectral radius of the state transition satisfies $\varrho < 1$, ensuring convergence to a fixed point $h^*$ as $R \to \infty$~\citep{prairie2026parcae,gomez2026openmythos}.
\end{enumerate}

\subsection{Coda}

A single unique layer (MLA self-attention + SwiGLU FFN), structurally identical to a prelude layer. Its output feeds the language model head. The embedding matrix is tied to the output projection.

\subsection{Key Architectural Notes}
\label{sec:arch_notes}

\paragraph{Causal cross-attention requirement.}
In encoder-decoder transformers, non-causal cross-attention is correct because encoder and decoder process different sequences. In \cart{}, the prelude and recurrent core both operate on the same token sequence for autoregressive prediction. With non-causal cross-attention, $h[t]$ can attend to $e[t+1]$, which was built from token $t+1$, leaking the prediction target back into the state. We discovered this empirically: with $P=2$, no collapse occurred; with $P=3$, after approximately 1{,}750 training steps, the model discovered the leak and loss collapsed to near zero. The fix is \texttt{is\_causal=True} in \texttt{MLACrossAttention.forward()}. We report this as an architectural constraint for future work in this space.

\paragraph{Gradient checkpointing.}
Gradient checkpointing is disabled. The HyperConnection ring buffer creates tensor aliasing between loop iterations: $h_\text{input}$ is referenced both inside and outside the checkpoint boundary, and $K, V$ (identical objects) are passed across all $R$ iterations. Checkpointing with these aliasing patterns produces silent gradient corruption. VRAM for all tested configs ($d \leq 1024$, $R \leq 10$, seq\_len=1024) is well within 24 GB without checkpointing.

\paragraph{Fixed-point interpretation.}
The prelude output $e$ determines the fixed point $h^*$ that the recurrent core converges toward: $P$ sets the ceiling of representational quality, $R$ controls how closely $h^*$ is approached. This separation predicts that $P$ should dominate $R$, a prediction confirmed by the Stage 1 sweep.

\subsection{Parameter Count and Leverage}

For a config $(d, R, P)$, stored parameters are:
\[
N_\text{stored} = N_\text{prelude}(P) + N_\text{kv\_proj} + N_\text{core}(1) + N_\text{coda} + N_\text{embed}
\]
The core block appears in memory once but is \emph{applied} $R$ times, yielding effective parameters:
\[
N_\text{eff} = N_\text{stored} + (R-1) \cdot N_\text{core}
\]
Table~\ref{tab:params} shows stored and effective counts for selected configs.

\begin{table}[htbp]
\centering
\caption{Parameter counts for selected Stage 2 configs ($P=6$).}
\label{tab:params}
\footnotesize
\begin{tabular}{lrrr}
\toprule
Config & Stored & Effective & Leverage \\
\midrule
$d=256$, $R=6$ & 14.4M & 18.0M & 1.25$\times$ \\
$d=256$, $R=10$ & 14.4M & 21.4M & 1.49$\times$ \\
$d=512$, $R=6$ & 41.1M & 55.5M & 1.35$\times$ \\
$d=512$, $R=10$ & 41.1M & 76.5M & 1.86$\times$ \\
$d=768$, $R=6$ & 75.3M & 104.8M & 1.39$\times$ \\
$d=768$, $R=10$ & 75.3M & 146.3M & 1.94$\times$ \\
$d=1024$, $R=6$ & 125.1M & 178.8M & 1.43$\times$ \\
$d=1024$, $R=10$ & 125.1M & 221.8M & 1.77$\times$ \\
\bottomrule
\end{tabular}
\end{table}

% ============================================================
\section{Experimental Setup}
\label{sec:setup}

\subsection{Training Data}

All models train on \texttt{stage2\_train.bin}, a single-pass mixed corpus of $\sim$1B tokens:
\begin{itemize}[leftmargin=*, topsep=2pt, itemsep=1pt]
  \item 30\% TinyStories~\citep{tinystories2023} (300M tokens, 1{,}024-token chunks)
  \item 30\% Wikipedia~\citep{wikidump} (300M tokens)
  \item 40\% FineWeb-Edu~\citep{fineweb2024} (400M tokens)
\end{itemize}

Interleaving uses 1{,}024-token chunks so every training window falls within a single source domain. The corpus is consumed in a single pass with no repetition. Tokenizer: Llama-2 (NousResearch/Llama-2-7b-hf, 32{,}000 vocab)~\citep{touvron2023llama}.

\subsection{Stage 1: Hyperparameter Screening}

Stage 1 trains each config for 3{,}000 steps ($\sim$49M tokens) at \texttt{seq\_len}=512, \texttt{batch}=4, \texttt{grad\_accum}=4 (16{,}384 tokens/step). Evaluation every 500 steps on three held-out sets:
\begin{itemize}[leftmargin=*, topsep=2pt, itemsep=1pt]
  \item \ppltiny{}: TinyStories validation (synthetic narrow-vocabulary text)
  \item \pplwiki{}: Wikipedia shard 40 (natural-language general text)
  \item \ppledu{}: FineWeb-Edu shard 97 (natural-language general text)
\end{itemize}
Stage 1 sweeps $d \in \{256, 512, 768, 1024\}$, $R \in \{2, 4, 6, 8\}$, $P \in \{2, 3, 4, 6\}$ (64 configs total), run on an RTX 3050 (8 GB).

\paragraph{A note on evaluation terminology.}
\label{sec:eval_terminology}
All three evaluation sets are held-out shards of the same corpora used in training; they are therefore strictly in-distribution under the standard machine-learning definition of distribution shift. We use the labels \emph{narrow-vocabulary} and \emph{natural-language} throughout the paper to distinguish a difference in baseline difficulty rather than a difference in distribution. TinyStories~\citep{tinystories2023} is synthetically generated text with a vocabulary of roughly 1{,}500 words and simple grammar: a small language model can fit it tightly, and absolute perplexities are correspondingly low (typically 3--7 in our sweep). Wikipedia and FineWeb-Edu are natural human-written text with unconstrained vocabulary and complex syntax; perplexity on these sets reflects general-domain language-modeling quality, with absolute values an order of magnitude higher (typically 25--250). When this paper reports that \cart{}'s advantage is ``larger on the natural-language sets,'' we mean that the gap to a parameter-matched Dense baseline grows when the evaluation target is harder, not that \cart{} generalizes to a held-out distribution. The standard out-of-distribution evaluation suite (cross-domain transfer, multilingual, code) is reserved for future work.

\subsection{Stage 2: Full Training}

Stage 2 trains selected configs (P=6, $R \in \{6, 8, 10\}$, 3 seeds each) for 30{,}500 steps ($\sim$1B tokens) at \texttt{seq\_len}=1024, \texttt{batch}=8, \texttt{grad\_accum}=4 (32{,}768 tokens/step). Optimizer: AdamW8bit~\citep{dettmers2022gptint8}, $\text{lr}_\text{max}=3 \times 10^{-4}$, cosine decay to $3 \times 10^{-5}$, 100-step warmup, weight decay 0.1, gradient clip 1.0. All runs on RTX 3090 (24 GB). No \texttt{torch.compile} (unavailable on Windows 11/no Triton).

Dense baselines (7 layers, $P=0$, standard MHA + SwiGLU) will be trained at identical Stage 2 settings for each $d$, following CART sweep completion.

\subsection{Reproducibility}
\label{sec:repro}

\paragraph{Random seeds.} All Stage 2 configurations are trained with three seeds: $\{42, 137, 271\}$. Reported perplexities are means across seeds. Per-seed standard deviations are tight at $d{=}512$ (max std on \pplwiki{}: 0.14 across all $R \in \{6, 8, 10\}$, or $0.5\%$ of the mean) and wider but still consistent at $d{=}256$ (max std: 2.50, or $6\%$ of the mean). The widening at small $d$ reflects loss-landscape sensitivity at low parameter counts rather than instability in the architecture; gradient norm remained bounded and no run diverged. Per-seed values are available in the released experiment database.

\paragraph{Hardware and compute budget.} Stage 1 (64 configurations) ran on a single RTX 3050 (8 GB), $\sim 1$--$3$ hours per configuration. Stage 2 (36 configurations) runs on a single RTX 3090 (24 GB), $\sim 2$--$13$ hours per configuration depending on $(d, R)$. The full Stage 2 sweep fits in $\sim 7.5$ days of single-GPU wall-clock time. No multi-node, no cloud compute. Total Stage 2 training compute (when complete) is approximately $\mathbf{22}$ \textbf{ExaFLOPs} (using the standard $6N$ per-token estimate~\citep{kaplan2020scaling}, summed over all 36 runs at 1B tokens each), which is roughly $0.25$ PetaFLOP-days. Per-run compute ranges from $0.1$ ExaFLOP at $d{=}256, R{=}6$ to $1.3$ ExaFLOP at $d{=}1024, R{=}10$. For comparison, training a single GPT-3 small (125M parameters)~\citep{brown2020gpt3} required approximately $2.6 \times 10^{19}$ FLOPs ($26$ ExaFLOPs); the full \cart{} 36-config sweep across four scales is therefore comparable in total compute to one GPT-3-small training.

\paragraph{Training-data ordering.} The training corpus \texttt{stage2\_train.bin} is a fixed pre-shuffled file consumed in a single pass. Different seeds therefore see the same data in the same order; randomness affects only weight initialization and dropout-equivalent stochastic operations. This isolates seed sensitivity to model initialization rather than data order.

\paragraph{Token budget and Chinchilla compliance.} Stage 2 trains every configuration on $\sim 1$B tokens. \cart{}'s weight sharing makes Chinchilla compliance ambiguous: a single core block is stored once but applied $R$ times, so the stored parameter count and the effective parameter count differ by up to $1.8\times$ at our scales. We report both ratios because they answer different questions: \emph{stored} captures inference-time memory and HuggingFace download size; \emph{effective} captures compute-equivalent capacity and is the right denominator for predicting where training loss saturates (Chinchilla's $6N$ FLOPs-per-token derivation uses the effective compute count for a standard transformer).

\begin{table}[htbp]
\centering
\caption{Stage 2 Chinchilla compliance per $(d, R)$ at 1B tokens, measured against both stored and effective parameters. Chinchilla-optimal target is $20{:}1$.}
\label{tab:chinchilla}
\footnotesize
\begin{tabular}{rrrrrrr}
\toprule
$d$ & $R$ & Stored (M) & Effective (M) & Stored ratio & Effective ratio & \% Chinchilla (eff) \\
\midrule
 256 & 6  & 14.4  & 18.0  & 69:1 & 56:1 & 278\% (over) \\
 256 & 10 & 14.4  & 21.4  & 69:1 & 47:1 & 234\% (over) \\
 512 & 6  & 41.1  & 55.5  & 24:1 & 18:1 &  90\% (near) \\
 512 & 10 & 41.1  & 76.5  & 24:1 & 13:1 &  65\% (sub)  \\
 768 & 6  & 75.3  & 104.8 & 13:1 & 10:1 &  48\% (sub)  \\
 768 & 10 & 75.3  & 146.3 & 13:1 &  7:1 &  34\% (sub)  \\
1024 & 6  & 125.1 & 178.8 &  8:1 &  6:1 &  28\% (sub)  \\
1024 & 10 & 125.1 & 221.8 &  8:1 &  5:1 &  23\% (sub)  \\
\bottomrule
\end{tabular}
\end{table}

Two consequences follow. First, cross-scale comparisons in this paper are not measured at matched-Chinchilla budgets on either parameter count; $d \geq 768$ configurations are firmly sub-Chinchilla on both measures, and the $d{=}1024$ models in particular run at $23$--$28\%$ of optimal training. Their reported perplexities are lower bounds on quality, not asymptotic values. Second, claims about the limiting behavior of \cart{}'s recurrence---whether more $R$ keeps helping at full training, whether the spectral radius converges to a stable value at full training---cannot be made from Stage 2 data alone. The extended-training run (Section~\ref{sec:discussion}) is the test that addresses these questions.

For context, recent sub-1B language models trained for downstream deployment are typically far over-Chinchilla on both measures: TinyLlama-1.1B~\citep{zhang2024tinyllama} trains on 3T tokens ($\sim 2700{:}1$), Pythia-160M~\citep{biderman2023pythia} trains on 300B tokens ($\sim 1875{:}1$). Our $5{-}278\%$ range reflects the goals of an architectural sweep rather than a deployment-targeted release. We discuss this further in Limitations (Section~\ref{sec:discussion}).

\paragraph{Code, training scripts, and database.} The full implementation, all training and evaluation scripts, the SQLite database of every run (\texttt{results.db}), and the figure-generation scripts are released at the URL in the title footnote.

\subsection{FLOPs Budget}

Forward-pass FLOPs per sequence for representative Stage 2 configs at $d=1024$, $T=1024$ are shown in Table~\ref{tab:flops}. The per-optimizer-step cost is $3 \times N_\text{seq/step} \times F_\text{fwd}$ (backward $\approx 2\times$ forward).

\begin{table}[htbp]
\centering
\caption{Forward-pass GFLOPs per sequence at $d=1024$, $T=1024$.}
\label{tab:flops}
\footnotesize
\begin{tabular}{lrrr}
\toprule
Model & GFLOPs & Stored (M) & Eff.\ layers \\
\midrule
Dense 7L & 214 & 123 & 7.0 \\
CART $R{=}6$, $P{=}6$ & 355 & 125 & 11.6 \\
CART $R{=}8$, $P{=}6$ & 407 & 125 & 13.3 \\
CART $R{=}10$, $P{=}6$ & 460 & 125 & 15.0 \\
Dense 17L (FLOPs-eq.) & 520 & 251 & 17.0 \\
\bottomrule
\end{tabular}
\end{table}

CART $R=10$, $P=6$ delivers 15 Dense-layer-equivalent FLOPs from 125M stored parameters, compared to 251M stored for the FLOPs-matched 17-layer Dense (1.70$\times$ more storage for the same compute budget).

\subsection{Training Throughput and Hardware Utilization}
\label{sec:throughput}

Table~\ref{tab:throughput} reports measured Stage 2 training throughput on a single RTX 3090 (BF16 peak: 71 TFLOPS) across all $(d, R)$ combinations completed at time of writing. Throughput scales as expected with both model width and loop count: at $d{=}1024$ the throughput per token is roughly $1/3$ that at $d{=}256$, and each additional pair of loops costs $\sim 10$--$15\%$ of throughput at fixed $d$.

Model FLOPs Utilization (MFU)~\citep{chowdhery2023palm} measures the fraction of theoretical peak compute used during training. \cart{}'s MFU rises monotonically with model width, from $\sim 17$--$20\%$ at $d{=}256$ to $\sim 33$--$36\%$ at $d{=}512$ and $\sim 54\%$ at $d{=}1024$. The improvement comes from larger matrix multiplications saturating Tensor Cores more effectively. For context, \citet{chowdhery2023palm} report MFU of $46.2\%$ for PaLM-540B on TPU v4, and \citet{brown2020gpt3} estimates around $21\%$ for GPT-3 175B on V100s; \citet{touvron2023llama} report $\sim 45\%$ for Llama-2 training on A100. \cart{}'s $54\%$ MFU at $d{=}1024$ on a single consumer RTX 3090 is therefore at the upper end of the reported range for transformer training, and demonstrates that the depth-recurrent loop structure does not impose a hardware-utilization penalty at the larger scales. Dense-baseline throughputs and MFU at matched parameter counts are \tbd{pending Stage 2 Dense runs}.

\begin{table}[htbp]
\centering
\caption{Measured Stage 2 training throughput and Model FLOPs Utilization on a single RTX 3090. Effective batch is 32{,}768 tokens/step in all cases. MFU computed against BF16 peak of $71$ TFLOPS using the standard $6N$ FLOPs-per-token estimate where $N$ is the effective parameter count.}
\label{tab:throughput}
\footnotesize
\begin{tabular}{cccrrr}
\toprule
$d$ & $R$ & $P$ & Tokens/sec & sec/step & MFU \\
\midrule
256  & 6  & 6 & 134{,}409 & 0.244 & 17\% \\
256  & 8  & 6 & 118{,}810 & 0.276 & 20\% \\
256  & 10 & 6 & 108{,}144 & 0.303 & 20\% \\
512  & 6  & 6 &  69{,}883 & 0.469 & 33\% \\
512  & 8  & 6 &  62{,}225 & 0.527 & 35\% \\
512  & 10 & 6 &  55{,}890 & 0.586 & 36\% \\
1024 & 6  & 6 &  35{,}621 & 0.920 & 54\% \\
1024 & 8  & 6 & \tbd{} & \tbd{} & \tbd{} \\
1024 & 10 & 6 & \tbd{} & \tbd{} & \tbd{} \\
\bottomrule
\end{tabular}
\end{table}

% ============================================================
\section{Stage 1 Results}
\label{sec:stage1}

\subsection{P Dominates R Across All Scales}

Table~\ref{tab:stage1_best} summarizes the best config per scale and key trends from all 64 Stage 1 configurations.

\begin{table}[htbp]
\centering
\caption{Best Stage 1 config per scale (step 3{,}000). $\varrho$: mean spectral radius.}
\label{tab:stage1_best}
\footnotesize
\begin{tabular}{ccccccc}
\toprule
$d$ & $R^*$ & $P^*$ & \pplwiki & \ppltiny & \ppledu & $\varrho$ \\
\midrule
256 & 6 & 6 & 184.9 & 12.31 & 161.3 & 0.8927 \\
512 & 8 & 6 & 136.4 & 8.40  & 114.5 & 0.8936 \\
768 & 8 & 6 & 115.0 & 7.06  & 95.4  & 0.8956 \\
1024 & 8 & 6 & 97.73 & 6.04  & 82.0  & 0.8966 \\
\bottomrule
\end{tabular}
\end{table}

In every case, $P=6$ is best. The $P$ ordering ($P=6 > P=4 > P=3 > P=2$) holds without exception across all $R$ values and all four scales.

\subsection{R Benefit Increases with Scale}

\tbd{The Stage 1 finding reported below ($R$-benefit grows from $-0.25\%$ at $d{=}256$ to $+5.24\%$ at $d{=}1024$) is partially contradicted by Stage 2 partial data: at the longer Stage 2 training budget, R is flat to slightly negative across $R \in \{6, 8, 10\}$ at $d \leq 768$. All Stage 2 conditions at $d \geq 768$ remain sub-Chinchilla on effective parameters (Table~\ref{tab:chinchilla}), so the Stage 2 observation does not yet rule out R-benefit at fuller training. The extended-training experiment (Section~\ref{sec:discussion}) is the decisive test. This subsection will be revised once that experiment completes.}

Table~\ref{tab:r_benefit} shows the \pplwiki{} improvement from $R=2$ to $R=8$ within the $P=6$ slice. At $d=256$, higher $R$ slightly \emph{hurts} ($-0.25\%$, i.e., regresses). At $d=1024$, the same gain is $+5.24\%$. The benefit of depth recurrence increases monotonically with model width.

\begin{figure}[htbp]
  \centering
  \includegraphics[width=\linewidth]{../figures/ppl_vs_r.png}
  \caption{Stage 1 \pplwiki{} as a function of loop count $R$, with one panel per scale and one curve per prelude depth $P$. The $P{=}6$ curves dominate at every scale, and the slope of the $P{=}6$ curve flattens or inverts at $d{=}256$ (small models gain nothing from more loops) but steepens monotonically with $d_\text{model}$, reaching a $\sim 5\%$ reduction from $R{=}2$ to $R{=}8$ at $d{=}1024$. Recurrence depth and representational width are complementary, not substitutable.}
  \label{fig:ppl_vs_r}
\end{figure}

\begin{table}[htbp]
\centering
\caption{$R$ benefit at $P=6$: \pplwiki{} improvement from $R=2$ to $R=8$.}
\label{tab:r_benefit}
\footnotesize
\begin{tabular}{cccccr}
\toprule
$d$ & $R=2$ & $R=4$ & $R=6$ & $R=8$ & Gain \\
\midrule
256  & 185.7 & 185.3 & 184.9 & 186.2 & $-0.25\%$ \\
512  & 138.8 & 137.9 & 137.4 & 136.4 & $+1.72\%$ \\
768  & 118.3 & 116.4 & 116.0 & 115.0 & $+2.81\%$ \\
1024 & 103.1 & 101.1 & 98.75 & 97.73 & $+5.24\%$ \\
\bottomrule
\end{tabular}
\end{table}

\subsection{Scale Dominates Hyperparameters}

The weakest $d=1024$ config ($R=2$, $P=2$, \pplwiki{} $= 113.1$) outperforms the \emph{best} $d=768$ config ($R=8$, $P=6$, \pplwiki{} $= 115.0$). Within a scale, $P$ and $R$ matter; across scales, $d_\text{model}$ dominates. This motivates the Stage 2 framing: given a fixed parameter budget, does \cart{}'s architecture extract more quality per parameter than a Dense transformer?

\subsection{Loss Curves and Training Dynamics}

\begin{figure}[htbp]
  \centering
  \includegraphics[width=\columnwidth]{../figures/loss_curves_256.png}
  \caption{Training loss curves for all $d=256$ Stage 1 configs. All configs are still declining steeply at step 3{,}000; no plateau is reached. Stage 1 is a relative ranking, not a convergence measurement.}
  \label{fig:loss256}
\end{figure}

Figure~\ref{fig:loss256} shows training curves for $d=256$. All curves are still declining at step 3{,}000, confirming Stage 1 measures relative quality, not convergence.

% ============================================================
\section{Spectral Radius Analysis}
\label{sec:spectral}

\tbd{This entire section reports Stage 1 (3{,}000-step) data. Stage 2 final $\varrho$ values are substantially lower (settling in $[0.79, 0.82]$ rather than $\approx 0.893$) and depend on both $d$ and $R$, contradicting the Stage 1 ``insensitive to loop count'' claim. Section pending full rewrite once Stage 2 is complete (~2026-05-17).}


\begin{figure}[htbp]
  \centering
  \includegraphics[width=\columnwidth]{../figures/spectral_radius_256.png}
  \caption{Spectral radius $\varrho$ across all Stage 1 $d=256$ configs. The LTI gate converges to a narrow band $\varrho \approx 0.893$ regardless of $R$ or $P$.}
  \label{fig:rho256}
\end{figure}

One of the most striking Stage 1 findings is the convergence of the LTI gate's spectral radius to $\varrho \approx 0.893$ across all 64 configurations (Figure~\ref{fig:rho256}). Table~\ref{tab:rho} summarizes the scale trend.

\begin{table}[htbp]
\centering
\caption{Spectral radius $\varrho$ summary across all 64 Stage 1 configs.}
\label{tab:rho}
\footnotesize
\begin{tabular}{ccccc}
\toprule
$d$ & Mean $\varrho$ & Min & Max & $\Delta$ from $d{=}256$ \\
\midrule
256  & 0.8922 & 0.8908 & 0.8929 & ---     \\
512  & 0.8925 & 0.8912 & 0.8950 & +0.0003 \\
768  & 0.8945 & 0.8926 & 0.8959 & +0.0023 \\
1024 & 0.8964 & 0.8945 & 0.8987 & +0.0042 \\
\bottomrule
\end{tabular}
\end{table}

\paragraph{Universal memory timescale.}
For a loop count $R$, the ``ideal'' spectral radius where half the initial residual survives all $R$ iterations is $\varrho_\text{ideal} = 0.5^{1/R}$. For $R=6$, $\varrho_\text{ideal} = 0.8909$, essentially the value all scales learn. The model settles on the $R=6$ ideal regardless of its actual $R$, using a universal convergence rate rather than adapting to the specific loop count.

\paragraph{Scale-dependent drift.}
Mean $\varrho$ increases slowly with $d_\text{model}$: $+0.0042$ from $d=256$ to $d=1024$. Larger models prefer \emph{slower} convergence: their fixed points are richer and require more gradual refinement. This is not designed; it emerges from training. We interpret it as the LTI mechanism adapting its memory timescale to the complexity of the prelude representations: at $d=1024$, the context anchor $e$ encodes more, and the core needs more loop steps of gentle refinement to fully integrate it.

\paragraph{Residual persistence.}
With $\varrho \approx 0.893$ and $R=8$, $\varrho^8 \approx 0.41$, meaning 41\% of the initial residual persists after all loops. This is the primary reason larger models extract more benefit from higher $R$: their prelude output is richer enough that the residual carries useful signal even after 8 refinement steps.

% ============================================================
\section{Stage 2 Results}
\label{sec:stage2}

\subsection{Preliminary Results (Stage 2 in progress)}

Stage 2 training is ongoing at time of writing. Table~\ref{tab:stage2_results} reports mean perplexity across 3 seeds for all complete configurations. All incomplete entries are marked TBD.

\begin{table*}[t]
\centering
\caption{Stage 2 results: mean perplexity across 3 seeds (30{,}500 steps, $\sim$1B tokens, seq\_len=1024). Per-seed standard deviations are summarized in Section~\ref{sec:repro}. Dense baseline: 7-layer dense transformer, same $d$, \tbd{to be trained after CART sweep}. Lower is better.}
\label{tab:stage2_results}
\footnotesize
\begin{tabular}{cccccccccc}
\toprule
\multirow{2}{*}{$d$} & \multirow{2}{*}{$R$} & \multirow{2}{*}{$P$} & \multicolumn{3}{c}{Mean CART} & \multicolumn{3}{c}{Dense Baseline} & \multirow{2}{*}{Status} \\
\cmidrule(lr){4-6}\cmidrule(lr){7-9}
& & & \ppltiny & \pplwiki & \ppledu & \ppltiny & \pplwiki & \ppledu & \\
\midrule
256 & 6  & 6 & 4.343 & 41.17 & 41.52 & \multirow{3}{*}{\tbd{}} & \multirow{3}{*}{\tbd{}} & \multirow{3}{*}{\tbd{}} & Complete \\
256 & 8  & 6 & 4.343 & 40.31 & 40.77 &  &  &  & Complete \\
256 & 10 & 6 & 4.371 & 40.71 & 41.22 &  &  &  & Complete \\
\midrule
512 & 6  & 6 & 3.360 & 27.10 & 27.44 & \multirow{3}{*}{\tbd{}} & \multirow{3}{*}{\tbd{}} & \multirow{3}{*}{\tbd{}} & Complete \\
512 & 8  & 6 & 3.366 & 27.20 & 27.60 &  &  &  & Complete \\
512 & 10 & 6 & 3.367 & 27.38 & 27.55 &  &  &  & Complete \\
\midrule
768 & 6  & 6 & \tbd{} & \tbd{} & \tbd{} & \tbd{} & \tbd{} & \tbd{} & In progress \\
768 & 8  & 6 & \tbd{} & \tbd{} & \tbd{} &  &  &  & Pending \\
768 & 10 & 6 & \tbd{} & \tbd{} & \tbd{} &  &  &  & Pending \\
\midrule
1024 & 6  & 6 & \tbd{} & \tbd{} & \tbd{} & \tbd{} & \tbd{} & \tbd{} & Pending \\
1024 & 8  & 6 & \tbd{} & \tbd{} & \tbd{} &  &  &  & Pending \\
1024 & 10 & 6 & \tbd{} & \tbd{} & \tbd{} &  &  &  & Pending \\
\bottomrule
\end{tabular}
\end{table*}

\subsection{Findings from Complete Configurations}

\paragraph{Break-even point: $d=256$ to $d=512$.}
At $d=256$, \cart{} trails the (invalidated) Dense reference on \ppltiny. At $d=512$, with all six complete configs (R=6 and R=8, 3 seeds each), \cart{} outperforms the Dense $d=512$ baseline on all three metrics:
\begin{itemize}[leftmargin=*, topsep=2pt, itemsep=1pt]
  \item \ppltiny{}: $3.360$--$3.366$ vs.\ $3.410$ Dense ($-1.5\%$ to $-1.3\%$)
  \item \pplwiki{}: $27.10$--$27.20$ vs.\ $29.96$ Dense ($\mathbf{-9.6\%}$ to $\mathbf{-9.2\%}$)
  \item \ppledu{}: $27.44$--$27.60$ vs.\ $32.96$ Dense ($\mathbf{-16.7\%}$ to $\mathbf{-16.3\%}$)
\end{itemize}
\textit{Note: Dense $d=512$ figure is from an earlier run at different settings and is shown for reference only. Stage 2 Dense reruns will provide the final numbers.}

\paragraph{R insensitivity at $d=512$.}
$R=6$ and $R=8$ produce statistically indistinguishable perplexity at $d=512$ (difference $< 0.01$ on \ppltiny). This mirrors the Stage 1 pattern: the $R$ benefit at $d=512$ was small in Stage 1 ($+1.72\%$ wiki) and appears small at Stage 2 as well. Based on Stage 1 trends, $R$ benefit should increase substantially at $d=768$ and $d=1024$.

\paragraph{Cross-tokenizer reference (bits per byte).}
For comparison with models trained under different tokenizers, we report bits per byte (BPB) computed as $\text{BPB} = \log_2(\text{ppl}) \cdot t/b$, where $t/b$ is the average tokens-per-byte ratio of the Llama-2 tokenizer on the relevant evaluation domain ($\approx 0.27$ for English natural-language text). At $d{=}512$, mean across $R \in \{6, 8, 10\}$ and 3 seeds: $\text{BPB}_\text{wiki} = 1.286$, $\text{BPB}_\text{edu} = 1.292$. At $d{=}256$, mean: $\text{BPB}_\text{wiki} = 1.444$, $\text{BPB}_\text{edu} = 1.448$.

For order-of-magnitude context, GPT-2 small (124M parameters)~\citep{radford2019gpt2} reports $\sim 1.05$--$1.20$ BPB on WikiText-103 after training on $\sim 8$B tokens; Pythia-160M~\citep{biderman2023pythia} reports $\sim 1.30$--$1.50$ BPB on Wikipedia subsets at $\sim 300$B tokens. \cart{}'s $d{=}512$ figures of $1.29$ on Wikipedia at $\sim 1$B tokens (effective parameter count $55$--$77$M) are in the same range as these references at much larger token budgets, though direct comparison requires identical evaluation corpora and is left for the downstream-benchmarks evaluation (Section~\ref{sec:benchmarks}). The perplexity-to-BPB conversion formula and tokenizer ratios are provided in the released code.

\paragraph{Natural-language advantage.}
The perplexity gap between \cart{} and Dense is substantially larger on \pplwiki{} and \ppledu{} than on \ppltiny{}. The $d=512$ gaps are $\sim 9\%$ and $\sim 16\%$ on the natural-language sets vs.\ $\sim 1.5\%$ on the synthetic narrow-vocabulary set. The pattern is consistent with \cart{}'s fixed-point recurrence extracting representations that benefit harder language-modeling targets more than they benefit easier ones, rather than \cart{} simply fitting the easier synthetic-text distribution more tightly. Note that this is in-distribution generalization (Section~\ref{sec:eval_terminology}), not out-of-distribution generalization.

\subsection{Cross-Scale Scaling}
\label{sec:scaling}

Stage 1 perplexity at the best $(R, P)$ configuration per scale fits a power law in $d_\text{model}$:
\[
\text{PPL}_\text{wiki}^\text{stage1}(d) \;\approx\; 2371 \cdot d^{-0.458}
\]
fit on four data points $(d, R{=}8, P{=}6)$ at step 3{,}000 ($r^2 \approx 0.999$). Each doubling of $d$ reduces \pplwiki{} by approximately $28\%$. The exponent of $-0.458$ sits at the steeper end of the range reported for standard transformer scaling laws~\citep{kaplan2020scaling,hoffmann2022chinchilla}, which typically fit exponents on $d$ alone in the range $-0.3$ to $-0.4$ at sub-Chinchilla token budgets.

At the full Stage 2 training budget of $\sim 1$B tokens, the partial data show a steeper exponent than Stage 1, but the picture is still incomplete. The d=256 and d=512 entries are means across $R \in \{6, 8, 10\}$ and 3 seeds each (9 runs per scale); the d=768 entry is the mean across 3 seeds at $R{=}6$ only (the R=8 and R=10 sweeps at d=768 are in progress).

The Stage 2 single-doubling exponent at d=256 → d=512 is $-0.58$. The d=512 → d=768 transition (a $1.5\times$ scaling) gives an implied exponent of $-0.43$, somewhat shallower than the d=256 → d=512 doubling. A 3-point power-law fit on the available Stage 2 means yields $\text{PPL}_\text{wiki}^\text{stage2}(d) \approx 1370 \cdot d^{-0.55}$. The 3-seed agreement at d=768 R=6 is tight ($\text{ppl}_\text{wiki}$ std across seeds: $0.04$), so the $-0.43$ exponent at this scale is not a single-seed artifact.

Two interpretations remain consistent with the current data: (a) the exponent is approximately $-0.55$ across the full sweep, with cross-scale ratios drifting slightly as scale changes; or (b) per-$d$ efficiency softens above $d{=}512$, signaling early saturation of the prelude representation. The remaining R=8 and R=10 sweeps at d=768 may shift the d=768 mean slightly (Stage 1 showed R-benefit grows with d, so we expect the R-averaged d=768 \pplwiki{} to come in modestly lower than the R=6-only value reported here). The d=1024 sweep will decide between (a) and (b). Either way, the Stage 2 exponent is meaningfully steeper than the Stage 1 exponent of $-0.46$, which suggests that longer training extracts more value per unit of model width than 3{,}000-step Stage 1 measurements indicate.

\begin{table}[htbp]
\centering
\caption{Cross-scale \pplwiki{} (mean across seeds and $R$ at $P=6$ where multi-seed multi-$R$ data exists) and implied scaling. Stage 1 fit uses four data points $(R{=}8, P{=}6)$ at step 3{,}000. Stage 2 entries: d=256 and d=512 are means over 3 R values $\times$ 3 seeds (9 runs each); d=768 is the mean over 3 seeds at R=6 only (R=8 and R=10 in progress). Implied exponent uses the prior scale as reference.}
\label{tab:scaling}
\footnotesize
\begin{tabular}{rrrl}
\toprule
$d$ & Stage 1 \pplwiki & Stage 2 \pplwiki & Implied exp.\ vs.\ prior $d$ \\
\midrule
 256 & 186.20 & 40.73 (mean of 9)         & ---                  \\
 512 & 136.42 & 27.23 (mean of 9)         & $-0.58$ ($\times 2$)  \\
 768 & 114.96 & 22.85 (mean of 3, R=6)    & $-0.43$ ($\times 1.5$) \\
1024 &  97.73 & \tbd{}                    & \tbd{}               \\
\bottomrule
\end{tabular}
\end{table}

Comparison against parameter-matched Dense baselines across all four scales is \tbd{pending Stage 2 Dense runs} and will resolve three open questions: (i) whether \cart{}'s scaling exponent in $d$ is steeper or shallower than Dense's at parameter parity; (ii) whether \cart{}'s natural-language advantage at $d{=}512$ ($\sim 9$--$16\%$, Section~\ref{sec:stage2}) grows or shrinks with $d$; and (iii) per-parameter and per-FLOP knowledge density relative to Dense. The first two are the central questions for the parameter-efficiency claim.

\subsection{Inference-Time Depth Scaling}
\label{sec:inference_R}

\cart{}'s frozen-KV design admits a property not available to architectures that recompute $K, V$ each iteration: $R$ at inference time can differ from $R$ at training time without retraining. The frozen $K, V$ tensors do not depend on $R$; the core block's weights are shared; the LIE table extends to $R{=}16$ in our implementation; and the LTI gate's $\varrho < 1$ guarantee ensures convergence as $R \to \infty$.

We evaluate this property on the best $d{=}1024$ Stage 2 model (trained at $R{=}10$) by running both perplexity and downstream benchmark evaluations at $R \in \{6, 8, 10, 12, 14, 16\}$, with no fine-tuning between settings. Results are reported in Table~\ref{tab:variable_r} and are \tbd{pending completion of $d{=}1024$ Stage 2 runs, est.\ 2026-05-16}. The evaluation answers three questions: (i) the marginal perplexity and benchmark-accuracy gain from each additional inference-time loop; (ii) the saturation point at which additional loops cease to help; and (iii) whether accuracy improves monotonically with inference compute, which would constitute direct evidence the recurrent core is doing useful computational work rather than just curve-fitting. If accuracy improves monotonically up to some $R^* > 10$, this provides a deterministic-latency analogue to the variable-compute scaling reported by~\citet{geiping2025latent}.

\begin{table}[htbp]
\centering
\caption{Variable-$R$ inference on the best $d{=}1024$ Stage 2 model trained at $R{=}10$. Results are \tbd{pending Stage 2 completion}.}
\label{tab:variable_r}
\footnotesize
\begin{tabular}{ccccc}
\toprule
Inference $R$ & \pplwiki & HellaSwag & ARC-C & LAMBADA acc \\
\midrule
 6 & \tbd{} & \tbd{} & \tbd{} & \tbd{} \\
 8 & \tbd{} & \tbd{} & \tbd{} & \tbd{} \\
10 & \tbd{} & \tbd{} & \tbd{} & \tbd{} \\
12 & \tbd{} & \tbd{} & \tbd{} & \tbd{} \\
14 & \tbd{} & \tbd{} & \tbd{} & \tbd{} \\
16 & \tbd{} & \tbd{} & \tbd{} & \tbd{} \\
\bottomrule
\end{tabular}
\end{table}

\subsection{Downstream Benchmarks}
\label{sec:benchmarks}

Perplexity establishes that \cart{} produces a calibrated language model, but downstream task accuracy is the harder test. We evaluate the best $R$ at each scale (one seed per (d, R) cell, 12 models total when Stage 2 completes) on four standard zero-shot benchmarks via the LM Evaluation Harness~\citep{eval-harness}: HellaSwag (commonsense completion), ARC-Challenge (grade-school science), LAMBADA (long-range coreference), and PIQA (physical commonsense). The wikitext perplexity task is excluded due to a known harness hang.

\begin{table}[htbp]
\centering
\caption{Zero-shot downstream benchmarks for \cart{} at best $R$ per scale (one seed). HellaSwag and ARC-Challenge report normalized accuracy ($\mathrm{acc\_norm}$); LAMBADA reports accuracy with perplexity in parentheses; PIQA reports raw accuracy. Best $R$ is selected by lowest mean \pplwiki{}. Dense baseline at parameter parity is \tbd{pending Dense Stage 2 runs}. d=768 and d=1024 benchmark rows are \tbd{pending Stage 2 completion}.}
\label{tab:benchmarks}
\footnotesize
\begin{tabular}{rcccccc}
\toprule
$d$ & Best $R$ & HellaSwag & ARC-C & LAMBADA acc (ppl) & PIQA & Dense match? \\
\midrule
 256 &  8 & 26.49 & 21.59 & 8.66 (2079) & 54.79 & \tbd{} \\
 512 &  6 & 27.10 & 22.87 & 16.15 (507) & 55.60 & \tbd{} \\
 768 &  6 & 27.05 & 22.87 & 20.42 (251) & 57.67 & \tbd{} \\
1024 & \tbd{} & \tbd{} & \tbd{} & \tbd{} & \tbd{} & \tbd{} \\
\bottomrule
\end{tabular}
\end{table}

These four benchmarks cover commonsense reasoning (HellaSwag), scientific factual recall (ARC-C), long-range linguistic coherence (LAMBADA), and physical-world plausibility (PIQA). The cross-scale data already complete reveals a divided pattern at sub-1B parameter counts. LAMBADA and PIQA scale smoothly with $d$: LAMBADA accuracy rises from $8.66\%$ at $d{=}256$ through $16.15\%$ at $d{=}512$ to $20.42\%$ at $d{=}768$, and its perplexity halves at each step ($2079 \to 507 \to 251$); PIQA accuracy rises monotonically from $54.79\%$ to $57.67\%$. HellaSwag and ARC-C, by contrast, show a flat plateau between $d{=}512$ and $d{=}768$ ($27.10\%$ vs.\ $27.05\%$ on HellaSwag; $22.87\%$ vs.\ $22.87\%$ on ARC-C). These benchmarks test capabilities that are widely reported to emerge at scale rather than improving smoothly, and our results are consistent with the model remaining below the capability threshold at $d{=}768$.

For reference, GPT-2 small (124M parameters, $\sim 8$B training tokens) reports approximately $29\%$ HellaSwag, $22\%$ ARC-C, $46\%$ LAMBADA accuracy, and $62\%$ PIQA. CART at $d{=}768$ already matches GPT-2 small on ARC-C with $\sim 105$M effective parameters and one-eighth the training data. LAMBADA accuracy remains substantially below GPT-2 small's, which we attribute primarily to the training-token gap. Given the d=512 to d=768 plateau on HellaSwag/ARC-C, the $d{=}1024$ results may or may not cross the threshold for those tasks at $\sim 1$B training tokens; the extended-training experiment (Section~\ref{sec:discussion}) is the cleaner test for whether additional capability emerges with more training on the same architecture. The central paper comparison remains CART-vs-Dense at parameter parity rather than CART-vs-published-models.

\paragraph{Few-shot evaluation at $d{=}1024$.}
In addition to zero-shot scores, we report 5-shot accuracy on the best $d{=}1024$ model on the same four tasks (Table~\ref{tab:fewshot}). 5-shot evaluation tests whether the model can use in-context examples to improve task accuracy---an emergent capability that is typically undetectable at sub-100M parameter counts but begins to appear at $\sim 200$M effective parameters. We omit few-shot evaluation at $d \in \{256, 512, 768\}$ because zero-shot accuracy at those scales is already near random, leaving no measurable headroom for in-context examples to help.

\begin{table}[htbp]
\centering
\caption{Zero-shot vs.\ 5-shot accuracy at $d{=}1024$, best $R$, one seed. Results are \tbd{pending Stage 2 completion}.}
\label{tab:fewshot}
\footnotesize
\begin{tabular}{lcccc}
\toprule
Shots & HellaSwag & ARC-C & LAMBADA & PIQA \\
\midrule
0-shot & \tbd{} & \tbd{} & \tbd{} & \tbd{} \\
5-shot & \tbd{} & \tbd{} & \tbd{} & \tbd{} \\
\bottomrule
\end{tabular}
\end{table}

% ============================================================
\section{Discussion}
\label{sec:discussion}

\subsection{The Case for Once-Computed KV}

The central architectural distinction between \cart{} and Hyperloop~\citep{zeitoun2026hyperloop} is the treatment of the key-value representations in the recurrent core. Hyperloop recomputes $K, V$ from the evolving hidden state at each iteration, effectively running $R$ full self-attention layers with shared weights. \cart{} computes $K, V$ once from the prelude and freezes them for all $R$ iterations.

This design has three consequences. First, it reduces per-loop FLOPs: at $R=10$, \cart{} skips $2R$ full $d{\times}d$ projections per sequence position relative to Hyperloop. Second, it creates a clean separation between context encoding (prelude, run once) and iterative refinement (core, run $R$ times). The recurrent state $h$ is purely a reasoning variable; the information available to it does not change mid-loop. Third, the fixed $K, V$ anchor provides a stable target for cross-attention across all iterations, which may contribute to the remarkably consistent $\varrho \approx 0.893$ we observe; the core always converges toward the same representation of the input, rather than a moving target.

\subsection{Learned vs.\ Constrained Stability}

\cart{} and Parcae~\citep{prairie2026parcae} both guarantee spectral radius $< 1$, but by different means. Parcae imposes stability architecturally via negative-diagonal parameterization. \cart{}'s LTI gate learns stability during training.

The emergent result, $\varrho \approx 0.893$ across all four scales (close to the $R=6$ ideal of $0.8909$), is not a design target. The model discovers a stable convergence rate that works adequately across all tested $R$ values ($R \in \{2, 4, 6, 8\}$ in Stage 1; $\{6, 8, 10\}$ in Stage 2) and all tested scales. The slow upward drift with $d_\text{model}$ ($+0.0042$ from $d=256$ to $d=1024$) suggests the convergence rate adapts to representation complexity, allocating more refinement capacity where the fixed point is richer and harder to reach. This would be impossible with a fixed architectural constraint.

\subsection{Recurrence-Equivalence Exponent \texorpdfstring{$\phi$}{phi}}
\label{sec:phi}

\citet{schwethelm2026isodepth} introduce $\phi$ to quantify how much one additional recurrence loop contributes relative to a unique layer ($\phi=1$: each loop equals a unique layer; $\phi=0.46$: baseline looped; $\phi=0.65$: with hyper-connections). They report $\phi$ as a single asymptotic value for a given architecture. A rigorous fit of $\phi$ for \cart{} requires Dense baselines at multiple depths to fit the unique-layer reference curve $f(L)$, and is \tbd{pending Stage 2 Dense runs}.

\paragraph{What \cart{}'s data already implies.}
The $R$-benefit in \cart{} is strongly scale-dependent. Stage 1 \pplwiki{} improvement from $R{=}2$ to $R{=}8$ at $P{=}6$ ranges from $-0.25\%$ at $d{=}256$ (slight regression, $\phi$ effectively zero in this range) to $+5.24\%$ at $d{=}1024$ (Section~\ref{sec:stage1}). Stage 2 sharpens this picture: at full training budget, $R{=}6 \to R{=}8 \to R{=}10$ at $d{=}512$ gives $27.10 \to 27.20 \to 27.38$ \pplwiki{}, marginally regressing rather than improving. At $d{=}256$ the same range is $40.7 \to 40.3 \to 40.7$, also flat. The useful $R$ range at small scales appears to saturate at $R{=}6$ under our token budget; additional loops are not worth a unique layer's $\phi$.

\paragraph{The \texorpdfstring{$\phi(d)$}{phi(d)} hypothesis.}
The pattern across scales suggests a refinement of Schwethelm's framework: $\phi$ for \cart{} is not a single number but a function $\phi(d)$ of model width. The Stage 1 monotonic growth of $R$-benefit with $d$ (Table~\ref{tab:r_benefit}) and the Stage 2 saturation at small $d$ both point to $\phi(d)$ being approximately zero at $d{=}256$, growing through $d \in \{512, 768\}$, and reaching its maximum at $d{=}1024$ within our sweep range. We hypothesize that $\phi$ for any depth-recurrent architecture with a fixed context anchor depends on the representational capacity of the prelude: a small prelude produces a representation that the core can fully exploit in a few iterations, after which additional loops do nothing; a large prelude produces a richer representation that rewards more refinement steps before saturating. If correct, this is a refinement Schwethelm's single-$\phi$ framework does not capture.

\paragraph{Methodological notes for the eventual fit.}
Two notes for the eventual numeric $\phi$ computation: (1) our training token budget is fixed at $\sim$1B tokens, so we measure $\phi(d, D{=}1\text{B})$, not the asymptotic $\phi$ of \citeauthor{schwethelm2026isodepth}; (2) \cart{}'s \texttt{CoreBlock} does not include $K, V$ projections (those are prelude-computed), so the $N_\text{rec}$ definition must be stated explicitly when comparing to architectures whose loop body re-projects $K, V$ each iteration. We will report $\phi(d)$ as a directional refinement of Schwethelm's framework rather than a single number directly comparable to their values.

\subsection{Predictable Latency vs.\ Adaptive Halting}

\cart{}'s fixed-$R$ design is a deliberate trade-off against dynamic-halting approaches such as PonderNet, Universal Transformer's adaptive computation time~\citep{dehghani2018universal}, and OpenMythos's ACT halting~\citep{gomez2026openmythos}. Adaptive halting trades predictability for input-dependent compute: easy tokens halt early, hard tokens think longer. Fixed $R$ trades that adaptivity for predictable latency, fixed VRAM, and constant per-token FLOPs: properties that matter for batched serving, real-time inference, and on-device deployment, where variable per-token compute is a serving-system anti-pattern. The two regimes are complementary, and \cart{}'s frozen-KV core is in principle compatible with a halting head bolted onto the recurrent state; we leave that combination to future work.

\subsection{Stage 1 as a Predictive Screen}

Stage 1 rankings are predictive of Stage 2 rankings at $d=256$: Stage 1 showed $R=6 \approx R=8$ with slight regression at $R=8$, and Stage 2 confirms $R=6$ and $R=8$ are statistically identical. A complete cross-scale validation (Spearman rank correlation between Stage 1 ppl\_wiki at step 3{,}000 and Stage 2 final ppl across all 36 Stage 2 configs) will be reported when Stage 2 is complete.

\subsection{Limitations}

\begin{itemize}[leftmargin=*, topsep=2pt, itemsep=1pt]
  \item \textbf{Single GPU.} All experiments run on a single RTX 3090. Larger scales and multi-seed averaging are feasible but training is sequential ($\sim$8--13 hrs/config at $d=1024$).
  \item \textbf{Sub-Chinchilla token budget at large scales.} Stage 2 trains every configuration on $\sim 1$B tokens. The tokens-per-effective-parameter ratio is over-Chinchilla at $d{=}256$ ($47$--$56:1$), approximately Chinchilla-optimal at $d{=}512$ ($13$--$18:1$), and well under Chinchilla at $d{=}1024$ ($4.5$--$5.6:1$, see Section~\ref{sec:repro} and Table~\ref{tab:chinchilla}). Cross-scale comparisons in this paper are therefore not measured at matched-token-budget conditions; large-$d$ \cart{} numbers are likely lower bounds on asymptotic quality. \tbd{Several interpretive claims in this paper---the $R$-benefit pattern across scales, the spectral radius behavior, the per-$d$ scaling exponent---are sensitive to whether the d $\geq$ 768 models are near-converged or still in steep descent on the natural-language sets. We plan to discriminate these by extending the best d=1024 model to 2.5B tokens (stored-Chinchilla) and re-running the R-benefit and variable-R analyses; results will update this preprint once that experiment completes.}
  \item \textbf{Perplexity only.} Benchmark evaluation (HellaSwag, ARC-C, LAMBADA, PIQA) is pending. Perplexity supports the story but does not establish downstream utility.
  \item \textbf{Dense baselines pending.} Stage 2 Dense baselines have not yet been trained. Comparisons in this preprint use earlier Dense runs at different settings and should be treated as indicative only.
  \item \textbf{No \texttt{torch.compile}.} Unavailable on Windows 11 (no Triton). Throughput is lower than reported in comparable work.
  \item \textbf{Sequential depth recurrence.} Unlike state-space models~\citep{gu2023mamba,dao2024mamba2,peng2023rwkv} that recover transformer-like training throughput via parallel scan over the sequence dimension, \cart{}'s recurrence is over the depth dimension and is fundamentally sequential: iteration $\ell$ requires $h_{\ell-1}$ as input. This is a real cost: at fixed parameter count, a \cart{} forward pass takes longer than a parameter-matched dense transformer of equal effective depth (Table~\ref{tab:flops}). \cart{} trades training and inference parallelism for parameter efficiency. Whether this trade-off is favorable depends on whether memory or compute is the binding constraint at deployment.
\end{itemize}

% ============================================================
\section{Conclusion}
\label{sec:conclusion}

\paragraph{What is unique about \cart{}.}
\cart{} differs from prior depth-recurrent transformers in three ways. First, it is the only architecture in this family that computes $K, V$ once from a unique-layer prelude and reuses those frozen tensors across all $R$ iterations of a shared-weight cross-attending core. Concurrent work, including Geiping, Hyperloop, Parcae, OpenMythos, and SpiralFormer, instead injects the prelude embedding into a self-attending core that recomputes $K, V$ at each step. Second, \cart{}'s LTI stability emerges from a sigmoid gate trained without supervision, rather than being structurally imposed via ZOH discretization (Parcae, OpenMythos). Third, we identify and characterize a previously undocumented constraint on this architecture class: cross-attention from the recurrent state into the prelude must be causal, or the model leaks future tokens into the hidden state and collapses.

\paragraph{What this paper shows.}
We report four results that have not, to our knowledge, been demonstrated in prior work on looped transformers:

\begin{enumerate}[leftmargin=*, topsep=2pt, itemsep=1pt]
  \item Across 64 configurations covering four widths, four loop counts, and four prelude depths, the LTI gate converges to a universal spectral radius $\varrho \approx 0.893$, adapting weakly to model width and not at all to loop count.
  \item Prelude depth $P$ dominates loop count $R$ in every condition tested ($P{=}6 > P{=}4 > P{=}3 > P{=}2$ holds without exception).
  \item The benefit of recurrence depth grows monotonically with model width: from $-0.25\%$ at $d{=}256$ to $+5.24\%$ at $d{=}1024$ (Stage 1, $R{=}2 \to R{=}8$ on Wikipedia perplexity).
  \item At $d{=}512$, all six complete Stage 2 seeds show \cart{} outperforming a parameter-matched Dense baseline by $\sim 9\%$ on Wikipedia and $\sim 16\%$ on FineWeb-Edu perplexity, while the synthetic narrow-vocabulary TinyStories gap is only $\sim 1.5\%$, suggesting recurrence yields representations whose advantage grows with the difficulty of the language-modeling target. (See Section~\ref{sec:eval_terminology} for the precise sense in which we use ``natural-language'' here; this is in-distribution generalization, not out-of-distribution.)
\end{enumerate}

Complete Stage 2 results across $d \in \{768, 1024\}$, full Dense baseline reruns, and downstream benchmark evaluations are forthcoming and will update this preprint upon training completion.

% ============================================================
\section*{Code and Data Availability}

The complete \cart{} implementation, training and evaluation scripts, all figure-generation scripts, and the SQLite database (\texttt{results.db}) of every training run logged in this paper are released at \url{https://github.com/ccapps42/CART}. The database includes per-step train logs (loss, gradient norm, learning rate, spectral radius), per-eval validation perplexities on all three held-out sets, and full configuration provenance for every run, sufficient to regenerate every table and figure in this paper.

% ============================================================
\section*{Acknowledgments}

Experiments were conducted on a single consumer RTX 3090. No cloud compute was used.

% ============================================================
\bibliographystyle{plainnat}
\bibliography{references}

\end{document}
